\chapter{Future Work}
\label{ch:future-work}

This chapter outlines future directions for research and development 
based on the empirical findings of this thesis. 

Building upon the capabilities of the \tech{ln-history} platform demonstrated, 
we list numerous avenues for extending these in \autoref{sec:ln-history-project}.

Secondly, because the empirical data highlighted inefficiencies in the gossip protocol, 
\autoref{sec:ln-gossip-protocol} explores its ongoing development. 
It discusses the upcoming \term{Gossip v2} protocol updates, 
currently being drafted by \citeauthor{lightningdevelopers_bolts_2018}, 
as well as fundamentally new designs to share gossip messages,
aiming to resolve these existing propagation bottlenecks.

After that, \autoref{sec:the-role-of-ln} steps back to discuss the evolution of the Bitcoin scaling ecosystem.
The \gls{ln} was envisioned in the whitepaper \citetitle{poon_bitcoin_2016} as a decentralized, 
\gls{p2p} network capable of making Bitcoin a globally viable medium of exchange~\cite{poon_bitcoin_2016}. 
Nonetheless, as discussed in \autoref{ch:conclusion}, 
the network has matured with central hubs forming a backbone for payment routing.
To finish off, an outlook on emerging Layer-2 solutions that either complement the \gls{ln} or 
offer alternative approaches to scaling Bitcoin payments, is presented.


\section{\tech{ln-history} Project}
\label{sec:ln-history-project}

Because the \tech{ln-history} platform was designed with a modular architecture, 
it provides a foundation for extensive future development. 
These improvements can generally be categorized into expanding the platform's 
analytical capabilities and scaling its underlying infrastructure.

\paragraph{Analytical Extensions}

The dataset of \gls{ln} gossip messages and peer session metadata opens
many directions for future empirical research. 
The platform could be extended to track and analyse the following network statistics:

\begin{description}[leftmargin=1em, style=nextline]
    \item[Node Implementation Heuristics:] 
        By analysing the \term{feature-bits} within \code{node\_announcements}, 
        the platform could classify the network share of different Lightning node client software. 
        This could be achieved by natively integrating the scanning heuristics developed by 
        \citeauthor{myers_impscan_2025}~\cite{myers_impscan_2025} into the \service{gossip-processor}.
    
    \item[ISP and Cloud Centralization:] 
        Extending the peer address data model to resolve Autonomous System Numbers (ASNs) 
        and TLS certificates would allow researchers to quantify the network's reliance on major cloud providers, 
        similar to the analysis done by the \citetitle{mempooldevelopers_mempool_2026} project~\cite{mempooldevelopers_mempool_2026}.
    
    \item[Tor Network Dynamics:] 
        Deploying additional collector nodes operating only behind the Tor network 
        would enable comparative studies to measure whether privacy-preserving nodes 
        suffer from increased propagation lag or topological isolation compared to clear-net nodes.
    
    \item[Advanced Closure Analysis:] 
        Expanding the \code{channel\_closures} management (done by \service{chain-enricher}) 
        could provide deeper insights if channel closures are uni- or bilateral.
        Additionally, an analysis of spent mining fees for closures and 
        the final liquidity distribution of the closed channel could be made.
    
    \item[Negative Routing Fees:] 
        The database could be updated to track the addition of negative routing fees by the \citeauthor{lnddeveloper_lightningnetwork_2026}~\cite{lnddeveloper_lnd_2024}, 
        allowing researchers to analyse the effect onto the network and which nodes use them.
    
    \item[Routing Algorithm Simulations:] 
        The historical snapshots generated by the \service{ln-history-api} provide the perfect empirical foundation 
        for back-testing and evaluating the efficiency of (new) pathfinding algorithms on real-world topologies.
        Extending the research done by \citeauthor{saraswathi_exposition_2024} in~\cite{saraswathi_exposition_2024}.
    
    \item[Fee Market Correlations:] 
        Pending the recovery of the missing \code{channel\_update} data for the year 2024, 
        future research could investigate whether node operators dynamically adjust their routing fees
        in response to transaction fees in the Bitcoin blockchain.
\end{description}


\paragraph{Infrastructural Scaling}

Currently, the platform operates on a single virtual machine utilizing \tech{Docker Compose}. 
While this was a solid, resource-efficient decision for the scope of this thesis, 
the planned addition of new collector nodes and heavier analytical workloads will require more resources. 

To achieve true \term{high-availability} and manage growing data volumes, 
the infrastructure should be migrated to \tech{Kubernetes}~\cite{burns_borg_2016}. 
Adopting this industry standard for container orchestration would enable a new degree of scalability, 
allowing individual services such as the \service{gossip-processor} and \service{chain-enricher}
to be scaled horizontally across a distributed cluster based on real-time network traffic.


\section{Gossip Protocol}
\label{sec:ln-gossip-protocol}

This section examines two current developments for the \gls{ln} gossip protocol. 
First, we outline the \term{Gossip v2} architecture to support new channel types. 
Secondly, we discuss the potential integration of \term{set reconciliation technologies} 
to optimize bandwidth and propagation efficiency.


\paragraph{Gossip v2}
The activation of the Taproot soft fork introduced the cryptographic foundation for \term{Taproot channels}~\cite{wuille_bip341_2020}. 
Unlike \textit{legacy} channels that rely on visible 2-of-2 \gls{p2wsh} multisignature scripts, 
Taproot channels utilize Schnorr signatures~\cite{wuille_bip340_2020} and MuSig2 key aggregation~\cite{nick_musig2_2021}. 
This allows the channel funding transaction to appear identically to a standard, single-signature transaction on-chain, 
improving network privacy by making Lightning channels look like regular Bitcoin transactions.

The legacy gossip protocol v1 is structurally incompatible with these new cryptographic primitives.
The proposed \term{Gossip v2} (sometimes referred to as v1.5 or v1.75 during the development phases) 
introduces native support for Schnorr signatures. 
The specification draft has just been merged on May 4th 2026~\cite{mouton_extensionbolt_2023}.


\paragraph{Set Reconciliation and Minisketch}

As concluded in \autoref{ch:conclusion}, the current \gls{ln} gossip protocol suffers 
from bandwidth inefficiencies and incorrect synchronization. 
\citeauthor{lightningdevelopers_bolts_2018} have explored alternatives to tackle these challenges. 
The underlying problem of efficiently synchronizing a state across a distributed network 
is not unique to the \gls{ln} but a fundamental challenge in \gls{p2p} networks.

For example, the Bitcoin network layer has faced similar transaction relay bottlenecks and 
explored \tech{Erlay}~\cite{naumenko_erlay_2019} as a potential mitigation. 
Originally proposed to reduce node bandwidth by replacing continuous message flooding
with periodic set reconciliation using a library called \term{Minisketch}, 
its practical integration seems highly complex. 
As highlighted in a Bitcoin Core developer discussion~\cite{fleon_erlay_2026}, 
\tech{Erlay} trades bandwidth efficiency for increased propagation latency. 
Because set reconciliation requires additional network roundtrips to compute the state differences 
before actual data transmission, it is by design slower than \textit{naive fanout} flooding. 

Developers are actively debating whether the theoretical bandwidth savings justify 
the substantial increase in complexity and the downgrade in propagation speed~\cite{fleon_erlay_2026}. 
This ongoing debate shows that moving away from flood-based routing is a delicate 
architectural decision rather than a trivial protocol upgrade.

\begin{figure}[H]
    \centering
    \begin{tikzpicture}[
        >=Stealth, 
        timeline/.style={dashed, gray, thick},
        action/.style={font=\small, align=center},
        box/.style={draw, rounded corners, fill=gray!5, text width=2.8cm, align=center, font=\footnotesize, inner sep=5pt}
    ]
        % We have two
        \node[draw, circle, thick, minimum size=1.2cm] (NodeA) at (0, 0) {Node A};
        \node[draw, circle, thick, minimum size=1.2cm] (NodeB) at (7, 0) {Node B};

        % Initial Sets
        \node[box, above=0.3cm of NodeA] {Initial Set:\\ $\{tx_1, tx_2, tx_3\}$};
        \node[box, above=0.3cm of NodeB] {Initial Set:\\ $\{tx_2, tx_3, tx_4\}$};

        % Vertical Timelines
        \draw[timeline] (NodeA) -- (0, -7);
        \draw[timeline] (NodeB) -- (7, -7);

        % Step 1: Send Sketch
        \draw[->, thick] (0, -1.5) -- (7, -2) node[midway, above, action] {1. Send Sketch(Set A)};

        % Step 2: Compute Difference (Bob's side)
        \node[anchor=west, action, text width=3.5cm] at (7.2, -2.5) {
            2. Compute Difference\\ 
            Missing in B: $\{tx_1\}$\\
            Missing in A: $\{tx_4\}$
        };

        % Step 3: Send missing & Request
        \draw[<-, thick] (0, -4) -- (7, -3.5) node[midway, above, action] {3. Send $\{tx_4\}$, Request $\{tx_1\}$};

        % Step 4: Final response
        \draw[->, thick] (0, -5) -- (7, -5.5) node[midway, above, action] {4. Send $\{tx_1\}$};

        % Final Synced Sets
        \node[box, below=0.3cm] at (0, -7) {Reconciled Set:\\ $\{tx_1, tx_2, tx_3, tx_4\}$};
        \node[box, below=0.3cm] at (7, -7) {Reconciled Set:\\ $\{tx_1, tx_2, tx_3, tx_4\}$};

    \end{tikzpicture}
    \caption{
        A simplified sequence diagram of set reconciliation. 
        Instead of sending full transaction lists, nodes exchange compact mathematical sketches 
        to efficiently compute the symmetric difference and synchronize their states.
    }
    \label{fig:set-reconciliation}
\end{figure}

As illustrated in \autoref{fig:set-reconciliation}, set reconciliation allows two nodes 
to synchronize their data by exchanging compressed \term{sketches}. 
When comparing these sketches, nodes can calculate the symmetric difference 
between their states and request only the specific missing messages and therefore reduce bandwidth usage.

The application of this technology to the \gls{ln} has been considered since the 
network's early days, as seen by \citeauthor{russel_diyhplus_2018}'s initial 
mailing list proposals in \citeyear{russel_diyhplus_2018}~\cite{russel_diyhplus_2018}. 
The concept was later refined by \citeauthor{myers_magical_2022} in the \term{Magical Minisketch} proposal~\cite{myers_magical_2022}. 
Most recently, active research has resumed on the \textit{Delving Bitcoin} forum, 
where developers like \citeauthor{harvey-buschel_jharveyb_2026} are developing custom tools (namely \tech{gossip-observer}) 
to gather the empirical data necessary to design a Minisketch-based gossip architecture~\cite{harvey-buschel_jharveyb_2026}.


\section{The Role of Lightning in the Bitcoin Ecosystem}
\label{sec:the-role-of-ln}

We take a step back and look at Bitcoins goal to be classified as money.
For that, an asset must fulfil the following three functions. 
It must serve as a \textit{Medium of Exchange}, a \textit{Unit of Account}, and a \textit{Store of Value}~\cite{jevons_money_1875}.
The \gls{ln} as a scaling solution focuses on making Bitcoin a viable Medium of Exchange.
This last section will first present the current state of the research that focuses on the limitations of the \gls{ln}.
Further, it introduces new designs, like \gls{mpc} as well as neighbouring projects 
that present an alternative to the design of the \gls{ln}.


\paragraph{Limitations of the Bitcoin Lightning Network}

While the \gls{ln} facilitates significantly faster and cheaper payments than the 
Bitcoin base layer, its theoretical scalability is still bounded by its topology. 
The formalization of these constraints, achieved through the mathematical modelling of payment channel networks, 
has been pioneered by \citeauthor{pickhardt_mathematical_2026}~\cite{pickhardt_mathematical_2026}.
In a recent presentation, \citeauthor{pickhardt_mathematical_2026} demonstrated 
that the network's maximum throughput is constrained by the rate of infeasible payments~\cite{pickhardt_everything_2025}. 
His model reveals a trade-off between transaction size, 
routing reliability, and overall network throughput. 
For example, to maintain a \qty{90}{\percent} likelihood of feasibility for larger payments 
(up to \qty{922259}{\sat} within the professional routing network\footnote{Selected subset of \textit{central} nodes and channels in the \gls{ln}}), 
the expected bandwidth is limited at \num{70} payments per second~\cite{pickhardt_everything_2025}. 
This throughput is only one order of magnitude higher than the capacity of the underlying Bitcoin blockchain.


\paragraph{Multi Party Channels}

To overcome these structural limitations, \citeauthor{pickhardt_mathematical_2026} 
advocates for the deployment of \term{\glspl{mpc}} and channel factories. 
By allowing multiple participants to share a single on-chain \gls{utxo}.
\glspl{mpc} could scale transaction throughput by several orders of magnitude. 
The integration would require big changes to the current \gls{ln} protocol and its client implementations.
This limitation in addition to other challenges (infrastructure management, user-experience) indicates that 
further development to elevate \gls{ln} might even exist outside the current Lightning protocol itself. 


\paragraph{Emerging Second Layer Solutions}

Alternative second layer architectures have emerged to tackle liquidity and user-experience challenges.
In \citeyear{keceli_bitcoindev_2023}, \citeauthor{keceli_bitcoindev_2023} introduced \tech{Ark}~\cite{keceli_bitcoindev_2023}. 
This protocol utilizes \glspl{vtxo} to facilitate off-chain payments without requiring users to manage inbound liquidity,
which is a major friction point within the \gls{ln}. 
Multiple independent implementations of this protocol are currently in active development. 
Similarly, the \tech{Spark} protocol, developed by \citeauthor{lightspark_buildonspark_2026}~\cite{lightspark_buildonspark_2026}, 
leverages statechains alongside a shared-signing architecture to seamlessly transfer \gls{utxo} ownership off-chain.


\paragraph{Third Layer Solutions and Ecash}

Furthermore, the ecosystem is also expanding into third layer solutions that utilize the \gls{ln} as backbone. 
\tech{Cashu}, a free and open-source \term{Chaumian ecash protocol} built for Bitcoin, has gained traction. 
It utilizes the \gls{ln} as infrastructure to facilitate value exchange between \term{Mints}, 
which in turn issue \term{ecash tokens} to users. 
These tokens can be transacted offline and with perfect anonymity, 
utilizing the blind signature cryptography invented by \citeauthor{chaum_blind_1983} in \citeyear{chaum_blind_1983}~\cite{chaum_blind_1983}.


\paragraph{Outlook: Balancing Sovereignty and Convenience}

These emerging alternative protocols vary in their custody design.
Solutions like \tech{Ark}, \tech{Spark}, \tech{Cashu} claim to be self-custodial
because the user keeps unilateral exit rights to the Bitcoin base layer.
Nonetheless, their architectures require a \enquote{central} server, 
such as an Ark Service Provider (ASP), a Spark Service Provider (SSP), or a Cashu Mint
to coordinate liquidity or route payments. 

This mediated architecture stands in contrast to the \gls{ln} where
node operators maintain absolute control over their bitcoin. 
Ultimately, these developments indicate that the future of Bitcoin scaling 
might be a balance between sovereignty and convenience. 